\section{Continuous design, correlations, and convex performance}
\label{sec:a-design}

Fixed nominations give a different tractability boundary from uncertain
nominations.  On a cactus, a whole cycle can have arbitrarily many
correlated law parameters while its physical flow still has only one
free coordinate.  This section exploits that scalar equation.  It proves
exact local optimization and capacity feasibility, additive convex
minimization with rational original parameters, and a complementary
boundary for convex maximization.  The number and size of the cycles are
unrestricted throughout.

The optimization ingredients have direct predecessors.  Scalar threshold
bisection is classical quasiconvex optimization
\citep[Sections 2.1 and 3]{agrawal2020-dqcp}; the single-piece degree-one case is
linear-fractional optimization, and exact recovery from parametric linear
optimization goes back to \citet[Section 2]{megiddo1979-rational}.
The cycle capacity inequalities below already appear for quadratic laws
in \citet[Proposition 4.9, Lemma 4.10, and Proposition 4.11]
{amann2018-deciding-robust-feasibility-and-infeasibility}.
We give the arithmetic details needed to combine these mechanisms with
dense polynomial laws and to return rational admissible parameters,
including when an optimal physical flow is irrational.

\subsection{Exact recovery for a monotone polynomial root}

\begin{theorem}[Monotone root optimization over an H-polytope]
\label{thm:a-design-root}
Let $P=\{\theta\in\R^t:A\theta\le b\}$ be a nonempty bounded rational
polytope, with no bound on $t$.  On a supplied rational bracket $a<b$ let
\[
 H(q,\theta)=h_0(q)+\sum_{j=1}^t\theta_jh_j(q)
\]
be continuous and piecewise polynomial in $q$.  The partition consists
of fixed rational breakpoints independent of $\theta$, and every piece
has densely encoded rational coefficients and degree at most $D$.
Suppose, for every $\theta\in P$, that $H(\cdot,\theta)$ is strictly
increasing and $H(a,\theta)\le0\le H(b,\theta)$.
Then the smallest and largest roots over $P$ are computable exactly in
polynomial bit time in the total dense input.  Each output consists of a
rational optimizing vertex of $P$ and a real-algebraic root of degree at
most $D$ and polynomial coefficient bit length.  The theorem permits
lower-dimensional $P$, repeated polynomial roots, and roots at
breakpoints.  Strict increase and the bracket signs are promises.
\end{theorem}

\begin{proof}
Write $q(\theta)$ for the unique root.  At every rational $z\in[a,b]$,
strict increase gives the exact rational LP tests
\begin{align}
 \max_{\theta\in P}q(\theta)\ge z
 &\quad\Longleftrightarrow\quad
          \min_{\theta\in P}H(z,\theta)\le0,\label{eq:a-design-root-max}\\
 \min_{\theta\in P}q(\theta)\le z
 &\quad\Longleftrightarrow\quad
          \max_{\theta\in P}H(z,\theta)\ge0.\label{eq:a-design-root-min}
\end{align}
To justify both attainment and the notation for these extrema, express
any $\theta$ as a convex combination of vertices.  The same combination
of their $H(q(\theta),\cdot)$ values is zero.  At least one vertex has
nonpositive value and one has nonnegative value, so their roots bracket
$q(\theta)$.  The extreme vertex roots are therefore global extrema.
In particular the root map is quasilinear: each sublevel and superlevel
set is the intersection of $P$ with an affine halfspace.  No derivative
lower bound is needed.

We give one uniform separation bound over all vertex roots without
enumerating the vertices.  The case $t=0$ is ordinary single-polynomial
root isolation, with piece selection as described below; assume $t>0$.
Clear denominators in each row of $A\theta\le b$, and let $C\ge1$ bound
the absolute integer coefficients and right-hand sides.  Every vertex
has $t$ independent active rows: otherwise a nonzero direction
orthogonal to the active normals would admit small feasible
displacements in both directions.  This argument also applies to a
lower-dimensional polytope.  Cramer's rule gives a common nonzero
integer denominator $d$ and numerators $n_j$ with
\[
 |d|,|n_j|\le\Delta:=t!C^t.
\]
Clear every polynomial coefficient denominator, on every piece, by one
positive integer $R$, and let $P_0\ge1$ bound the resulting integer
coefficients.  Substituting a vertex and multiplying by $Rd$ gives an
integer polynomial of degree at most $D$ and height at most
\[
 H_0=(t+1)\Delta P_0.
\]
Delete duplicate breakpoints and zero-width pieces first.  Strict
increase implies that every specialized positive-width piece is
nonconstant, hence nonzero.  A breakpoint root belongs to an adjacent
closed piece by continuity.  Thus this height bound covers every
vertex root.  The logarithms of all these integers are polynomial in
the input length; taking products of input denominators only adds
their bit lengths.

Take any two of these nonzero integer polynomials, allowing equality,
and form their product $Q$.  It has degree at most $2D$ and height at
most $(D+1)H_0^2$.  Its primitive squarefree part $S$ is an integer
factor.  The classical factor bound of \citet{mignotte1974-factor}
gives $\|S\|_1\le2^{\deg S}\|Q\|_2$, so its height is at most
\[
 H_s=2^{2D}(2D+1)(D+1)H_0^2.
\]
For completeness, if $S$ has degree $n\ge2$, Cauchy's root bound puts
all its complex roots in $|z|\le1+H_s$.  Its discriminant is a nonzero
integer.  For any two distinct roots at distance $\delta$, the
discriminant product formula therefore implies
\[
 1\le H_s^{2n-2}\delta^2
            [2(1+H_s)]^{n(n-1)-2}.
\]
Since $n\le2D$, the rational number
\begin{equation}
 \sigma=[2(1+H_s)]^{-4D^2}
 \label{eq:a-design-separation}
\end{equation}
is a lower bound on every nonzero distance between roots of any two
vertex-piece polynomials.  If $n=1$ there is no pair to separate.
Squarefree extraction handles common factors and repeated roots.
Moreover $\log(1/\sigma)$ is polynomial in the dense input length.
The assumptions exclude $D=0$.

Bisect $[a,b]$ using \eqref{eq:a-design-root-max}, retaining $[l,u]$
containing the largest root, and stop when $u-l<\sigma$.  True tests,
including equality, raise $l$; false tests lower $u$.
Now minimize $H(l,\theta)$ over $P$ and select a vertex optimizer
$\theta_v$.  Its value is nonpositive, so
\[
 l\le q(\theta_v)\le q_{\max}\le u.
\]
Both middle quantities are vertex roots; separation forces equality.
A vertex optimizer is obtainable by imposing the rational optimal
objective equality and then minimizing coordinates successively on
the remaining face.  After $t$ further LPs all coordinates are fixed.
Each intermediate set is a face of the original $P$ and contains an
original vertex.  Hence each newly fixed coordinate has the original
Cramer bound, and the initial objective optimum is the objective at
an original vertex.  All intermediate LP encodings have a common
polynomial bound, rather than an iterated unspecified size bound.
For the minimum, bisect with \eqref{eq:a-design-root-min} and take a
maximizing LP vertex at the final upper bracket endpoint; its root
lies between $q_{\min}$ and that endpoint.  The same separation argument
applies.  Endpoints of the initial bracket cause no exception.

Finally substitute the recovered vertex, evaluate $H$ at the ordered
rational breakpoints and at $a,b$, and return a rational zero if one
occurs there.  Otherwise their signs locate a positive-width piece
with one root.  Its polynomial, squarefree part, and rational
isolating interval give the exact encoding.  Standard univariate
isolation and refinement have polynomial bit complexity in the degree,
coefficient bits, and requested accuracy bits
\citep{basu2006-algorithms-in-real-algebraic-geometry}.
All bisection queries have polynomial-bit rational coefficients.
Continuity makes a query at a breakpoint unambiguous on $P$.
\end{proof}

The theorem supplies an explicit exact recovery step for a classical
LP-bisection reduction.  It is not an algorithm for arbitrary roots of
a polynomial family or for a sparsely encoded binary exponent: the
precision in \eqref{eq:a-design-separation} depends polynomially on the
numerical degree.  Moving breakpoints and selection among nonmonotone
branches are also outside its scope.

\subsection{Cycle laws and capacities under correlated parameters}

Fix a balanced rational nomination $b$ on a connected simple cactus.
An admissible law has the form
\begin{equation}
 g_e(x;\theta)=g_{e0}(x)+\sum_j\theta_jg_{ej}(x),
 \label{eq:a-design-law}
\end{equation}
with fixed rational piece boundaries and dense rational polynomial
coefficients.  For every admissible parameter vector, the law is
promised continuous and strictly increasing on $\R$, with
$g_e(0;\theta)=0$.  The total piece count and degrees may grow.  These
promises are not an extension of the independent-box validation
algorithm in \cref{prop:a-law-validation} to general correlations.

Here is the scalar reduction, including its existence assumptions.
For a fixed parameter vector the primitive
$\int_0^x g_e(s;\theta)\,ds$ is nonnegative and strictly convex.
For $|x|\ge1$ it is at least $c_e(|x|-1)$, with the positive constant
\[
 c_e=\min\{g_e(1;\theta),-g_e(-1;\theta)\}.
\]
The sum is coercive, so it has a unique minimum on the nonempty
conservation affine space.  Stationarity supplies potentials and the
edge laws.  Each nonzero physical flow points to a lower potential;
the resulting directed graph is acyclic.  Source-to-sink flow
decomposition therefore bounds every magnitude by
$B:=\sum_v|b_v|$.

Conservation fixes every bridge flow and every cycle's effective
nominations independently of all law parameters.  Orient a cycle
consistently and choose its circulation $q_C$ to equal one reference
edge flow.  Rational conservation elimination gives oriented flows
$q_C+d_e$ with one $d_e=0$.  A reversed input edge uses
$-g_e(-x;\theta)$, which preserves all law promises, including for
nonodd laws.  Its consistency equation is
\begin{equation}
 H_C(q,\theta):=\sum_{e\in C}g_e(q+d_e;\theta)=0.
 \label{eq:a-design-cycle}
\end{equation}
It is strictly increasing in $q$, affine in $\theta$, and its root
lies in $[-B,B]$.  Its partition is the union of the translated edge
breakpoints.  Translation expands $(q+d_e)^j$ using only binomial
coefficients and powers of rational offsets, with polynomial bit cost
and degree unchanged.  Thus the union partition and every dense
coefficient list have polynomial size; no product of edge polynomials
is formed.  Conservation and these cycle equations are sufficient
for global potential consistency because the cactus cycles form a
basis of its cycle space.

\begin{theorem}[Global correlations: exact capacities and individual arcs]
\label{thm:a-design-global-capacity}
Let one parameter vector in \eqref{eq:a-design-law} range over a
nonempty bounded rational polytope $P$, allowing arbitrary correlations
between cycles.  Feasibility of rational signed arc bounds is one
rational linear feasibility problem.  A feasible instance has a
polynomial-time computable rational admissible parameter witness
satisfying all bounds exactly.  Robust satisfaction over $P$ is
decidable by rational LPs, which also return rational violating
scenarios when a bound fails.

Every individual signed arc has polynomial-time computable exact
algebraic extrema and rational optimizing parameter scenarios, also
when other such capacities filter the parameter domain.  The same
statements hold for rational linear flow inequalities whose
nonconstant part involves at most one cycle circulation.
\end{theorem}

\begin{proof}
Check bridge bounds directly.  In the original orientation a cycle
edge has flow $s_e(q_C+d_e)$, $s_e\in\{-1,1\}$.  Collect all its
bounds into a rational interval $a_C\le q_C\le b_C$, intersecting
with $[-B,B]$.  A permitted linear constraint is a constant test if
its circulation coefficient is zero, and otherwise another scalar
bound, with its inequality reversed if that coefficient is negative.
Reject any failed constant test or $a_C>b_C$.
For every fixed $\theta$, strict increase gives
\begin{equation}
 a_C\le q_C(\theta)\le b_C
 \quad\Longleftrightarrow\quad
 H_C(a_C,\theta)\le0\le H_C(b_C,\theta).
 \label{eq:a-design-capacity-lp}
\end{equation}
Both expressions are rational affine functions of $\theta$.  Conjoin
all these inequalities with $P$ to obtain a bounded rational polytope
$P_{\rm cap}$.  Rational linear feasibility gives precisely the claimed
decision and witness.  In particular irrational physical roots do not
require an algebraic parameter certificate.

For robust validation, maximize $H_C(a_C,\theta)$ and minimize
$H_C(b_C,\theta)$ over $P$ for each cycle.  All lower-bound maxima
must be nonpositive and all upper-bound minima nonnegative.  These
tests include equality; a strict failure supplies a rational LP
optimizer certifying a physical violation by
\eqref{eq:a-design-capacity-lp}.  A failed constant bound is violated
by any rational point of $P$.

Apply \cref{thm:a-design-root} to the target cycle over $P$, or over
the nonempty $P_{\rm cap}$ when capacities filter the scenarios.
Rational shifts and sign changes give the desired signed arc extrema
without increasing their algebraic degree.  Bridge extrema are fixed
rational values.  If $B=0$, all physical flows vanish: check the bounds
and return any rational parameter point when feasible.  This also
handles an edgeless graph before invoking a nontrivial root bracket.
\end{proof}

For quadratic laws, \eqref{eq:a-design-capacity-lp} is the
interval-to-halfspace reduction of A\ss mann and coauthors cited above;
conjoining it across cactus blocks preserves polyhedral uncertainty.
Affine polynomial coefficients give the same reduction by rational
evaluation.  Global correlations require no new convexity of the flow
image for this argument.  In particular this theorem makes no claim
that a globally correlated attainable flow region is a box or that
coupled convex performance can be minimized over that image.

\begin{theorem}[Independent cycles: attainable intervals and exact profiles]
\label{thm:a-design-independent}
Suppose instead that cycle $C$ has its own nonempty bounded rational
parameter polytope $P_C$, independent between cycles.  Bridge laws
are fixed admissible laws, or have independent nonempty rational
parameter polytopes of the same type.  Then the attainable flow region
is
\begin{equation}
 x=x^0+Zq,\qquad \ell_C\le q_C\le u_C,
 \label{eq:a-design-box}
\end{equation}
where $x^0$ is rational and the columns of $Z$ have disjoint cycle
supports and entries in $\{0,-1,1\}$.  Each endpoint has an exact
algebraic encoding of degree at most the largest law degree on its
cycle and a polynomial-time computable rational parameter profile.
Rational cycle-local capacities clip these intervals exactly and
preserve rational endpoint profiles, including singleton intervals;
infeasibility is decided exactly.  When feasible, every rational
linear flow objective has an exactly optimizing
rational parameter scenario, with or without these capacities.
Its scalar value is computable to any additive precision in polynomial
bit time; exact scalar comparison of sums across cycles is a separate
problem.
\end{theorem}

\begin{proof}
If $B=0$, every physical flow and reference circulation is zero.
The attainable intervals are all $[0,0]$, and capacities reduce to
constant tests at the zero flow.  Reject any failed test; otherwise
rational linear feasibility supplies an admissible point in each
cycle and variable bridge parameter polytope.  These rational profiles
attain every endpoint and optimize every linear flow objective, whose
value is rational.  This also handles the connected edgeless graph.
Henceforth $B>0$, so $[-B,B]$ is a nontrivial root bracket.

The root map of \eqref{eq:a-design-cycle} is continuous: strict signs
at points just below and above a root trap nearby roots, by joint
continuity in $(q,\theta_C)$.  This also proves one-sided continuity
at the common bracket endpoints.  The image of the compact connected
$P_C$ is therefore a compact interval.  Apply
\cref{thm:a-design-root} to get its endpoints and rational profiles.
Cycle independence and consistency give every point of
\eqref{eq:a-design-box}; bridges have fixed flows regardless of their
law parameters.

We will also need rational realization away from the endpoints.
A rational target $q$ is attainable exactly when the two rational LP
values satisfy
\[
 v^-:=\min_{\theta\in P_C}H_C(q,\theta)\le0
 \le v^+:=\max_{\theta\in P_C}H_C(q,\theta).
\]
Let $\theta^-,\theta^+$ attain them.  If $v^+>v^-$, set
\begin{equation}
 \lambda=\frac{-v^-}{v^+-v^-},\qquad
 \theta=(1-\lambda)\theta^-+\lambda\theta^+.
 \label{eq:a-design-interpolation}
\end{equation}
This rational profile belongs to $P_C$ and has polynomial bit length.
Its balance is exactly $H_C(q,\theta)=0$.
Affine constant terms are preserved by
convex combinations.  If $v^+=v^-$, both are zero and either profile
works.  No division by an edge flow is used.

Reduce the capacities to $[a_C,b_C]$ as in
\cref{thm:a-design-global-capacity}.  The constrained interval is
\[
 [L_C,U_C]=[\max\{\ell_C,a_C\},\min\{u_C,b_C\}].
\]
Exact comparison of the separately encoded algebraic endpoints and
rational clipping values decides emptiness in polynomial bit time.
Every new endpoint is either an old endpoint with its stored rational
profile, or a rational value realized by
\eqref{eq:a-design-interpolation}.  At ties either applicable witness
works.  A rational singleton is realized by interpolation; an
irrational singleton must inherit its endpoint from the original
interval and uses that endpoint's rational profile.  Combining
independent witnesses satisfies every capacity exactly.

A rational linear objective is a rational constant plus
$\sum_C\gamma_Cq_C$ with rational $\gamma_C$.  For maximization choose
$U_C$ when $\gamma_C>0$, $L_C$ when $\gamma_C<0$, and either when it
vanishes; minimization reverses the choices.  The resulting rational
profiles optimize exactly without comparing sums of algebraic values.
For additive value output, enclose each selected endpoint $q_C$ in a
rational interval $I_C$ and evaluate the objective at rational
representatives $\widehat q_C\in I_C$, for example their midpoints.
The rational estimate $\widehat v$ of the optimal value $v$ satisfies
\[
 |\widehat v-v|\le\sum_C|\gamma_C|\,
       \operatorname{width}(I_C).
\]
Stop when this sum is at most the requested error $\eps>0$; it
suffices to refine every interval to width at most
$\eps/(1+\sum_C|\gamma_C|)$.
This requires only polynomially many precision bits and rational
arithmetic operations.
The square-root-sum obstruction already present in the quadratic
special case of \cref{thm:a-weight-cactus-region} explains why these two
output contracts must be kept distinct.
\end{proof}

In particular, $g_e(x;\beta)=\beta_e x|x|$ with explicit positive
rational bounds and any rational polytope coupling the resistances
within one cycle satisfies the hypotheses with $D=2$.  Independent
positive resistance intervals are its box specialization.  These
specializations give exact capacitated continuous resistance design
and local quadratic algebraic extrema without a bound on cycle size
or parameter dimension.

\subsection{Convex minimization with exact capacity recovery}

\begin{lemma}[A rational Lipschitz extension from a cube]
\label{lem:a-design-perspective}
Let a rational polynomial $h$ be convex on $[-1,1]^k$, $k\ge1$,
with a polynomial-time exact rational value and gradient oracle.
Suppose rational bounds $V\ge|h|$ and
$G\ge\max_j|\partial_jh|$, $G\ge1$, are known on that cube.
Put
\begin{align*}
 M&=1+V+kG,\qquad t(z)=\max\{1,\|z\|_\infty\},\\
 \widehat h(z)&=t(z)h(z/t(z))+M(t(z)-1).
\end{align*}
Then $\widehat h$ is globally convex, agrees with $h$ on the cube,
and has an exact rational value and subgradient oracle.  It is
globally $kC$-Lipschitz in Euclidean norm for
$C=G+1+2(V+kG)$.  Oracle and bound encoding lengths are polynomial in
the input and query lengths whenever those of $h,V,G$ are.
\end{lemma}

\begin{proof}
Extend $h$ by $+\infty$ outside the cube.  The perspective
$s h(z/s)+M(s-1)$ is jointly convex on $s\ge1$, $|z_j|\le s$.
Perspective convexity and convexity under partial minimization are
standard \citep[Sections 3.2.5--3.2.6]{boyd2004-convex}.
For fixed $z$ and feasible $s$, set $w=z/s$.  Differentiation gives
\[
 \frac{\partial}{\partial s}\{s h(z/s)+M(s-1)\}
   =h(w)-\nabla h(w)^Tw+M\ge1.
\]
Thus its minimum over $s\ge t(z)$ occurs at $s=t(z)$ and equals
$\widehat h(z)$.  This proves global convexity and cube agreement.

To include boundary points and ties explicitly, choose any
$\tau\in\partial t(z)$ and put $w=z/t(z)$,
$K=h(w)-\nabla h(w)^Tw+M$.  The gradient of $h$ supports its
restriction to the cube even at its boundary.  Multiplying that
support inequality by the positive perspective variable gives
\[
 \widehat h(z')\ge\widehat h(z)+\nabla h(w)^T(z'-z)
                         +K(t(z')-t(z)).
\]
Here $1\le K\le1+2(V+kG)$, so the support inequality for $t$ yields
the subgradient $\nabla h(w)+K\tau$.  One may choose $\tau=0$ when
$\|z\|_\infty\le1$, and a signed maximizing coordinate vector
otherwise; ties permit any such vector.  Since $\|\tau\|_\infty\le1$,
this subgradient has infinity norm at most $C$.  Applying its support
inequality at both $z$ and $z'$ proves
\[
 |\widehat h(z)-\widehat h(z')|
       \le C\|z-z'\|_1\le kC\|z-z'\|_2.
\]
At rational queries, $t,w,K$ and the chosen subgradient use only
rational evaluation and comparisons of polynomial bit complexity.
\end{proof}

\begin{theorem}[Dense convex polynomial design]
\label{thm:a-design-convex}
In the independent-cycle model of \cref{thm:a-design-independent},
impose any of its permitted rational cycle-local capacities.  Let
$f\in\Q[x_1,\ldots,x_m]$ be densely encoded and promised convex on
$[-B-1,B+1]^m$, where $B=\sum_v|b_v|$.
For every rational $\eps>0$, one can decide infeasibility or return a
rational original parameter scenario satisfying every capacity
exactly and having objective at most $\eps$ above the constrained
minimum.  Time is polynomial in the total dense input and requested
accuracy bits, with no fixed bound on the number of cycles or the
polynomial degrees.  Convexity is a promise; sparse binary exponents,
variable nominations, potential bounds, and flow constraints coupling
distinct cycle coordinates are not covered.
\end{theorem}

\begin{proof}
First use \cref{thm:a-design-independent} to check feasibility and
store exact constrained endpoints $L_C,U_C$ with rational parameter
profiles.  If there are no edges, or $B=0$, check the fixed flow and
return any admissible rational profile.  Otherwise put $S=B+1$ and,
writing $f=\sum_\alpha c_\alpha x^\alpha$, compute
\[
 V_f=\sum_\alpha|c_\alpha|S^{|\alpha|},\qquad
 G_f=1+\max_i\sum_{\alpha:\alpha_i>0}
                    |c_\alpha|\alpha_i S^{|\alpha|-1}.
\]
These rational bounds have polynomial bit length under dense
encoding.  On the expanded flow box, $|f|\le V_f$ and
$\|\nabla f\|_\infty\le G_f$, whence
$|f(x)-f(y)|\le G_f\|x-y\|_1$ on every segment in that box.
Choose
\[
 \eta=\min\{1/2,\eps/(16mG_f)\}.
\]

Refine each constrained endpoint to rational enclosures
$L_C^-\le L_C\le L_C^+$ and $U_C^-\le U_C\le U_C^+$ of widths
at most $\eta$.  If $L_C^+\le U_C^-$, retain the rational inner
interval $[L_C^+,U_C^-]$.  Projection from the true interval moves a
coordinate by at most $\eta$.  Otherwise freeze the surrogate
coordinate at $L_C^+$ and retain the stored profile of the exact
endpoint $L_C$.  The true interval then has width less than
$2\eta$; every point in it lies within $2\eta$ of $L_C^+$, while
the stored endpoint is within $\eta$.  This works for an irrational
singleton even though it contains no rational target circulation.

Disjoint cycle supports imply that every feasible physical flow maps
to this rational surrogate box within $2m\eta$ in one-norm.
Conversely every surrogate flow becomes feasible by replacing only
its frozen coordinates with their stored exact endpoints, a change
of at most $m\eta$.  Every surrogate flow therefore lies in the
expanded box, since it differs from a feasible flow by at most
$\eta$ per edge.  The surrogate itself need not satisfy the physical
capacities on frozen cycles.

To optimize on the surrogate, remove every fixed coordinate,
including retained intervals of zero width, and affinely rescale the
remaining rational intervals to $[-1,1]^k$.  If $k=0$, evaluate the
single surrogate point.  Otherwise write its full flow as $x=Tz+a$
with rational $T,a$ and set $h(z)=f(Tz+a)$.  This is convex on the
cube.  Its exact rational value and gradient can be evaluated by
composition without expanding a new multivariate coefficient list.
Use $V=V_f$ and
$G=1+G_f\max_j\sum_e|T_{ej}|$ in
\cref{lem:a-design-perspective}.  Dense degrees bound rational powers
and their bit growth polynomially, including for rational oracle
queries outside the cube after the extension rescales them.

The resulting globally convex Lipschitz function has the explicit
rational bound $kC$ and a polynomial-time rational value oracle.
The cube contains a Euclidean unit ball and lies in a ball of radius
$k$, with an exact membership oracle.  The classical ellipsoid
optimization theorem, in the encoding formulation of
\citet[Theorem 2.5.9 and Section 2.5.1]{dadush2012-integer-programming},
therefore returns a rational point \emph{in} the cube within
$\eps/2$ of the minimum.  Its running time depends polynomially on
the lengths of the radii, Lipschitz bound, and accuracy.  Tiny
positive surrogate widths cause no exception: their rational affine
rescaling has polynomial bit length.  The extension agrees with
$h$ on the cube, so this is the required surrogate optimizer.

Realize every retained rational circulation exactly by
\eqref{eq:a-design-interpolation}, and use the stored endpoint
profile on each frozen cycle.  Choose any allowed rational bridge
parameters.  Every returned parameter is admissible, and its
physical flow satisfies all capacities exactly.  If $x^*$ minimizes
the true objective, project it to the surrogate for comparison with
the surrogate optimizer, then recover the latter as just described.
The total loss is at most
\[
 2mG_f\eta+\eps/2+mG_f\eta
       \le11\eps/16<\eps.
\]
All arithmetic used across cycles is rational.  Separate scalar root
enclosures also approximate the returned objective with the same
gradient bound; no common algebraic field or constitutive inverse
condition number is required.
\end{proof}

\begin{corollary}[Quadratic continuous resistance design]
\label{cor:a-design-quadratic}
For quadratic passive laws, independent positive rational resistance
intervals, or independent positive rational resistance polytopes
within cycles, admit the exact capacity and additive design guarantees
above for every rational objective
$f(x)=\tfrac12 x^TQx+d^Tx+c$ with $Q\succeq0$.
The conclusion also holds without capacities.
\end{corollary}
\begin{proof}
The piecewise quadratic laws satisfy \eqref{eq:a-design-law}, and the
objective is convex on the whole flow space.  Apply
\cref{thm:a-design-independent,thm:a-design-convex}.
\end{proof}

The significance of rational output is specific to design in the
original model.  An optimizer over rational approximations to an
algebraic circulation box need not itself be attainable or satisfy a
tight capacity.  Stored exact endpoint profiles and rational
interpolation resolve both issues.  Finite resistance sets lack this
interpolation property; the one-cycle realization hardness in
\cref{thm:a-bound-realization} precludes transferring this convex
minimization conclusion to finite choices.

\subsection{Convex maximization: coupling and measurement dimension}

\begin{theorem}[Bounded-data convex performance hardness]
\label{thm:a-design-maxcut}
Maximizing a rational convex quadratic function of all flows is
strongly $\NP$-hard on simple cacti of maximum degree three with a
fixed unit source and sink, all other nominations zero, and independent
resistance intervals contained in $[1,16]$.  All fixed resistances
also lie in that range.  Additive value approximation with absolute
error at most $1/4$ is $\NP$-hard.  On the explicit family below,
existential integer-threshold attainment is $\NP$-complete and robust
satisfaction of a half-integer upper bound is $\coNP$-complete.
\end{theorem}

\begin{proof}
Given an unweighted Max-Cut graph $J$ on $n$ vertices, construct a
chain of $n$ triangles linked by bridges.  Each triangle has distinct
entry and exit vertices: the direct edge joins them and the two-edge
alternate path uses its third vertex.  Attach bridges at entry and
exit, give every bridge resistance $4$, and nominate $+1$ at the
first entry and $-1$ at the last exit.  This is a simple cactus of
maximum degree three.  Orient direct edges, alternate paths, and
bridges from source to sink.  Each triangle carries unit through-flow.
Its alternate edge resistances are $2,2$ and its direct resistance
is an independent $\beta_i\in[1,16]$.  If $x_i$ is its direct flow,
\[
 \beta_i x_i^2=4(1-x_i)^2,
 \qquad x_i=\frac{2}{2+\sqrt{\beta_i}}\in[1/3,2/3].
\]
These coordinates vary independently, so $z_i=3x_i-1$ ranges over
the entire cube $[0,1]^n$.  All constructed physical flows are
positive.  The convex objective
\[
 F(x)=9\sum_{ij\in E(J)}(x_i-x_j)^2
          =\sum_{ij\in E(J)}(z_i-z_j)^2
\]
is a sum of squares of rational linear forms.  A convex function on
a cube has a maximizing vertex: fix coordinates successively and
choose an endpoint with value no smaller.  At binary $z$, $F$ is
exactly the cut size.  Consequently its maximum is the integer
maximum cut size of $J$.  This is the classical quadratic Max-Cut
encoding; see, for example,
\citet[Section 1.1]{delpia2017-miqp-np} for the binary formulation.

For integer $K$, existence of $F\ge K$ is precisely the Max-Cut
decision, and robust $F\le K-1/2$ is its complement.  Endpoint
choices $\beta_i=16$ for $z_i=0$ and $\beta_i=1$ for $z_i=1$ have
rational physical flows and give polynomial-size cut certificates.
This proves the stated restricted-family memberships, without an
NP membership claim for general algebraic performance values.
All physical numerical data are constants; expanded objective
coefficients and thresholds have polynomial magnitude in $|J|$.
The reduction remains polynomial under unary numerical encoding,
which proves strong hardness.  Rounding a value estimate of error
at most $1/4$ recovers the integer maximum and gives the fixed-error
claim.  This is an absolute error for the displayed objective,
without normalization to $[0,1]$.
\end{proof}

The physical network in this proof has only scalar cycle equations;
the comparison graph in the objective couples arbitrarily many
cycles.  Thus endpoint attainment does not make endpoint search
efficient.  A fixed number of measurements supplies the needed
restriction, through classical zonotope enumeration
\citep[Lemmas 2.1--2.3 and Theorem 2.6]{onn2004-convex}.

\begin{theorem}[A fixed number of linear measurements]
\label{thm:a-design-measurements}
Fix $k\ge1$.  In the independent polynomial-law model of
\cref{thm:a-design-independent}, including its permitted local
capacities, let $R\in\Q^{k\times m}$ and let $g$ be a densely encoded
rational polynomial convex on a box containing all attainable
measurements $Rx$.  One can decide infeasibility or return a
rational admissible parameter scenario
satisfying the capacities and maximizing $g(Rx)$ within additive
$\eps>0$.  This takes time $N^{O(k)}$ times a polynomial in
the requested accuracy bits, where $N$ is the dense input length.

The conclusion also holds for the fixed-nomination quadratic cactus
model with independent explicitly listed finite positive rational
resistance sets, without capacity filters.  No exact comparison of
all candidate values, or exactly maximizing scenario, is promised.
\end{theorem}

\begin{proof}
For the continuous model, first reject infeasible capacities using
\cref{thm:a-design-independent}; otherwise it supplies
the constrained affine box and a rational profile for each endpoint.
For finite quadratic sets, \cref{thm:a-weight-cactus-region} supplies
the same representation for the convex hull of the attainable flows,
and every box endpoint is attained by resistances belonging to the
original finite sets.  Denote the relevant bounds in either case
by $L_C,U_C$.  The measurement region, or its convex hull, is the
zonotope
\[
 y=Rx^0+\sum_C a_Cq_C,\qquad
 a_C=RZ_C\in\Q^k,\quad L_C\le q_C\le U_C.
\]
Its segment lengths can be algebraic, but its directions $a_C$ are
rational.  The convexity box contains the whole zonotope in the
finite case as well, since it contains its generating finite set.

Discard zero $a_C$ and enumerate the full-dimensional strict sign
cones of the central hyperplane arrangement $h^Ta_C=0$ in $\R^k$.
Here is an explicit rational implementation.  Insert the hyperplanes
one at a time.  For each current sign pattern try both signs on the
new hyperplane, retaining precisely those for which
\[
 \sigma_C h^Ta_C\ge1\quad\hbox{for every inserted }C
\]
is feasible.  Homogeneity makes this equivalent to strict sign
feasibility.  Rational LP gives a polynomial-bit representative.
The number of full-dimensional regions of $r$ hyperplanes in fixed
dimension is $r^{O(k)}$: adding a hyperplane creates at most the
number of regions of its induced arrangement in dimension $k-1$,
giving this bound by induction.  Repeated or dependent hyperplanes
cannot increase it.  Thus every insertion and the entire enumeration
use $N^{O(k)}$ bit operations up to polynomial rational LP factors.

For each retained sign pattern select $q_C=U_C$ when $h^Ta_C>0$
and $q_C=L_C$ otherwise, and combine the corresponding rational
endpoint profiles.  Zero $a_C$ can use either endpoint.  Every
vertex of the projected zonotope is obtained this way: a vertex
has a full-dimensional set of exposing directions in the ambient
space, even if the zonotope is lower dimensional.  Choose such a
direction off the finitely many nonzero-generator hyperplanes.
It uniquely selects the appropriate endpoint of every positive-length
segment.  Zero-length generators merely refine this arrangement and
do not change the selected point.  The point zonotope is included.
Convexity gives a maximum at a zonotope vertex.  In the finite case
that vertex is an original admissible scenario, so the finite and
convex-hull maxima agree.  Capacity filtering of finite sets is not
being inferred from this convex-hull argument.

It remains to compare candidates in bit complexity.  Put
\[
 T_0=B\max_j\sum_e|R_{je}|,\qquad
 M_0=\max_j\sum_C|a_{jC}|.
\]
Every true measurement lies in $[-T_0,T_0]^k$.  On
$[-T_0-1,T_0+1]^k$, bound every partial derivative of $g$ in
absolute value by a rational $L\ge1$ using the absolute monomial
coefficient sums, as in \cref{thm:a-design-convex}.  Convexity is
needed on the attainable zonotope, whereas this derivative bound
uses only polynomial evaluation and holds on the enlarged box
regardless of convexity there.  All bound encodings are polynomial
under dense degree encoding.

Refine every selected endpoint to rational error at most
\[
 \delta=\min\left\{\frac1{1+M_0},
               \frac{\eps}{8kL(1+M_0)}\right\}.
\]
Each measurement error is at most $M_0\delta$, stays in the
enlarged box, and changes the value by at most
$kLM_0\delta\le\eps/8$.  Dense univariate root refinement,
rational measurement sums, and dense evaluation of $g$ are
polynomial in input and accuracy bits.  Keep the candidate with
the largest rational estimated value and return its stored original
parameter scenario.  Two estimation errors cost at most $\eps/4$,
so this scenario has the claimed additive guarantee and retains
exact capacity feasibility in the continuous model.  This uses
no common field of all cycle roots and no exact comparison of
independent algebraic sums.

If $B=0$, all $a_C$ vanish, or $g$ is constant, choose any of the
stored feasible profiles.  The same construction covers zero-width
intervals and absence of cycles.  These cases avoid unnecessary
enumeration and do not affect the bound.
\end{proof}

\begin{corollary}[Fixed-rank quadratic maximization]
\label{cor:a-design-fixed-rank}
The preceding additive maximization algorithm applies to
$\tfrac12x^TQx+d^Tx+c$ when the rational positive-semidefinite
matrix $Q$ has fixed rank $r$, using at most $r+1$ measurements.
This holds in either scenario model of
\cref{thm:a-design-measurements}.
\end{corollary}
\begin{proof}
Rational $LDL^T$ elimination expresses the quadratic part as a sum
of $r$ nonnegative rational multiples of squares of rational linear
forms, with polynomial bit lengths.  A zero diagonal pivot in a
positive-semidefinite Schur complement has zero row and column and
can be omitted.  Use these forms and $d^Tx$ as measurements.  The
resulting polynomial in at most $r+1$ variables is convex.
\end{proof}

Fixed-rank convex binary quadratic maximization by zonotopes has the
direct predecessor \citet{ferrez2005-fixed-rank}; the general
fixed-dimensional edge-direction framework is due to
\citet{onn2004-convex}.  The application here supplies rational
directions, separately refined algebraic lengths, and admissible
rational passive-network scenarios.  These arithmetic and recovery
steps also extend the measurement theorem from quadratic resistances
to the independent polynomial-law polytopes proved above.  The
$N^{O(k)}$ bound is not a fixed-parameter tractability assertion in
$k$.  Together with \cref{thm:a-design-maxcut}, it isolates measurement
dimension as a meaningful restriction on coupled convex performance;
it does not transfer the finite-set maximization result to convex
minimization or exact target realization.
